\subsubsection{Kruskal}

\paragraph{Idea intuitiva.}
Encuentra el \textbf{arbol de expansion minima} (MST) en $O(E \log E)$. Ordena las aristas por peso ascendente y las aniade al MST si no forman ciclo (chequear con DSU). El arbol tiene exactamente $V - 1$ aristas. La demostracion (``Cut Property''): para cualquier corte, la arista minima cruzando el corte pertenece a \emph{algun} MST; por lo tanto, elegir la minima disponible que no forma ciclo es optimo.

\paragraph{Propiedades clave (invariantes).}
\begin{itemize}\setlength\itemsep{0pt}
	\item Algoritmo \textbf{greedy} correcto (Cut Property).
	\item Funciona en grafos \textbf{no conexos}; produce un \textbf{minimum spanning forest}.
	\item Si las aristas tienen pesos unicos, el MST es unico.
	\item El MST tiene exactamente $V - 1$ aristas (o menos si el grafo es disconexo).
\end{itemize}

\paragraph{Codigo clave.}
Kruskal con DSU:
\begin{verbatim}
sort(edges.begin(), edges.end());
DSU dsu(n);
long long mst = 0;
int edgesUsed = 0;
for (auto [w, u, v] : edges)
    if (dsu.unite(u, v)) {
        mst += w;
        if (++edgesUsed == n - 1) break;
    }
\end{verbatim}

\paragraph{Complejidad.}
\begin{itemize}\setlength\itemsep{0pt}
	\item $O(E \log E)$ (ordenamiento) o $O(E \log V)$ con DSU optimizations.
	\item DSU es casi $O(1)$ amortizado, asi que el sort domina.
\end{itemize}

\paragraph{Donde tocar para modificar.}
\begin{itemize}\setlength\itemsep{0pt}
	\item \textbf{Second best MST}: enumerar aristas del MST, sacar una, aniadir una no-MST; recalcular.
	\item \textbf{Kruskal randomizado}: si el grafo es casi-completo, agregar shuffle.
	\item \textbf{Aniadir info al DSU}: guardar el peso maximo en el camino, etc.
	\item \textbf{DSU con rollback}: si necesitas Kruskal reversible (para queries offline).
\end{itemize}

\paragraph{Trampas y errores frecuentes.}
\begin{itemize}\setlength\itemsep{0pt}
	\item El DSU debe estar limpio para cada invocacion.
	\item Si las aristas tienen pesos negativos, Kruskal sigue funcionando.
	\item Si el grafo es disconexo, el ``MST'' es un bosque con varias componentes.
	\item El check \texttt{if (++edgesUsed == n - 1) break;} optimiza la salida temprana.
\end{itemize}

\paragraph{Codigo.}
Ver \texttt{cpp.tex}, seccion \textit{Kruskal}.

\vspace{0.5em}
